\chapter{Hoarding Disorder}

\section*{Definition and Overview}
Hoarding disorder is a mental health condition recognised in diagnostic manuals as a distinct illness. It is characterised by a persistent difficulty discarding possessions, a strong urge to save items, and distress at the thought of parting with them, resulting in clutter that compromises the use of living spaces and causes significant distress or impairment. Unlike collecting or simple mess, hoarding disorder interferes with health, safety, relationships, and daily functioning. It is a chronic condition, but cognitive behavioural therapy specifically adapted for hoarding is effective.

\section*{Epidemiology}
Hoarding disorder affects approximately 2--6\% of the population, or more than half a million Australians. It is more common in older adults and often coexists with depression, anxiety disorders, and attention-deficit hyperactivity disorder. The disorder usually begins in early life but is often hidden for decades because people feel shame and avoid having others see their homes. It comes to attention most often through family concern, health and safety crises, or compulsory interventions. Both sexes are affected roughly equally.

\section*{Aetiology and Pathophysiology}
\begin{itemize}
    \item \textbf{Genetic vulnerability:} hoarding runs in families, with heritability estimated at around 50\%
    \item \textbf{Information processing:} deficits in categorising, organising, and decision-making
    \item \textbf{Cognitive beliefs:} exaggerated beliefs about the usefulness, beauty, or sentimental value of possessions
    \item \textbf{Emotional attachment:} objects are experienced as extensions of the self or as memory aids
    \item \textbf{Avoidance:} discarding triggers intense distress, so the person avoids it, reinforcing the behaviour
    \item \textbf{Trauma and loss:} onset and worsening are associated with adverse life events
    \item \textbf{Neurobiology:} altered activity in brain regions involved in decision-making and attachment
\end{itemize}

\section*{Clinical Features}
\begin{itemize}
    \item Difficulty discarding possessions even when they are broken, worthless, or hazardous
    \item Intense anxiety or distress when discarding is considered or attempted
    \item Clutter that prevents the use of rooms for their intended purpose
    \item Significant distress or impairment in social, occupational, or daily functioning
    \item Avoidance of visitors and social withdrawal
    \item Health and safety hazards: fire risk, falls, pests, and sanitation problems
    \item Hoarding of animals may occur in a minority, with serious welfare concerns
    \item The person may or may not recognise the problem (insight varies)
\end{itemize}

\section*{Differential Diagnosis}
\begin{itemize}
    \item Non-clinical clutter and collecting, which do not impair functioning
    \item Obsessive-compulsive disorder, where hoarding may be driven by obsessions and rituals
    \item Depression and apathy, which reduce motivation to maintain the home
    \item Dementia and cognitive decline, causing disorganisation
    \item Schizophrenia and delusional beliefs about possessions
    \item Autism spectrum conditions with special interests in objects
\end{itemize}

\section*{When to Seek Medical Care}
A GP is the first point of contact for assessment and referral. Families should seek help when they notice escalating clutter, health or safety concerns, or significant distress, rather than waiting for a crisis.

\textbf{Call 000 (or local emergency services) for:}
\begin{itemize}
    \item Fire or imminent fire risk in the home
    \item A vulnerable person unable to access food, medication, or care
    \item Severe self-neglect, immobility, or collapse
    \item Thoughts of self-harm or suicide
    \item Unsafe living conditions requiring urgent welfare intervention
\end{itemize}

\section*{Diagnosis and Investigations}
\begin{itemize}
    \item \textbf{Clinical interview:} the core features of persistent difficulty discarding, urges to save, distress, and clutter-related impairment
    \item \textbf{Standardised scales:} instruments such as the Saving Inventory measure severity
    \item \textbf{Mental health assessment:} exclude other disorders and identify coexisting conditions
    \item \textbf{Risk assessment:} fire, fall, health, and welfare risks
    \item \textbf{Home visit:} with consent, to assess clutter objectively
    \item \textbf{Capacity assessment:} where compulsory intervention is considered
\end{itemize}

\section*{Management}
\begin{itemize}
    \item \textbf{Cognitive behavioural therapy:} the first-line treatment, combining cognitive restructuring, decision-making skills, and graduated discarding
    \item \textbf{Home-based treatment:} sessions in the home give the best outcomes
    \item \textbf{Medication:} SSRIs and other antidepressants treat coexisting depression and anxiety, though no medicine is approved specifically for hoarding
    \item \textbf{Motivational interviewing:} builds engagement and readiness to change
    \item \textbf{Skills training:} organising, sorting, and maintaining gains
    \item \textbf{Multidisciplinary care:} psychologists, occupational therapists, and community services
    \item \textbf{Crisis and harm minimisation:} clearing hazards with consent, fire safety checks, and welfare coordination
    \item \textbf{Support groups:} peer support reduces shame and isolation
\end{itemize}

\section*{Complications}
\begin{itemize}
    \item Fire-related injury and death, falls, and trips in cluttered homes
    \item Hygiene problems, infestation, and respiratory illness
    \item Malnutrition and self-neglect
    \item Eviction, homelessness, and financial strain
    \item Relationship breakdown and social isolation
    \item Occupational impairment and disability
    \item Animal hoarding with cruelty and public health concerns
    \item Depression, anxiety, and suicide risk
\end{itemize}

\section*{Prognosis and Outcomes}
Hoarding disorder is a chronic condition, but treatment works. CBT tailored to hoarding leads to significant reductions in clutter and distress in the majority of people who engage, and improvements often persist with maintenance support. Outcomes are poorer when insight is limited, when treatment is imposed without consent, or when the person disengages. Early intervention, a collaborative approach, and treatment of coexisting conditions give the best results. Recovery is usually gradual, and progress is measured in small, sustained steps.

\section*{Prevention and Self-Care}
\begin{itemize}
    \item Seek help early and from practitioners experienced in hoarding
    \item Engage in home-based CBT and practise sorting and discarding skills
    \item Set small, achievable goals and celebrate progress
    \item Ask a trusted person for support with decisions about possessions
    \item Treat depression and anxiety alongside the hoarding
    \item Address fire and safety hazards first
    \item Attend support groups and stay connected socially
    \item Prepare for maintenance: relapse is common, and ongoing support helps
\end{itemize}

\section*{Sources and Further Reading}
The following reputable sources provide further information for healthcare students and patients seeking reliable, current advice about hoarding disorder.

\begin{itemize}
    \item Healthdirect. \textit{Hoarding disorder: symptoms and treatment.} https://www.healthdirect.gov.au/hoarding
    \item SANE Australia. \textit{Hoarding disorder support.} https://www.sane.org/
    \item Australian Psychological Society. \textit{Find a psychologist for hoarding.} https://psychology.org.au/
    \item NHS. \textit{Hoarding disorder: diagnosis and treatment.} https://www.nhs.uk/mental-health/conditions/hoarding-disorder/
    \item Mayo Clinic. \textit{Hoarding disorder: diagnosis and treatment.} https://www.mayoclinic.org/
    \item International OCD Foundation. \textit{Hoarding center.} https://hoarding.iocdf.org/
\end{itemize}
